\documentclass[11pt,a4paper]{article}
% preamble.tex — estilo común de todos los apuntes Froth.
% Cada apunte hace % preamble.tex — estilo común de todos los apuntes Froth.
% Cada apunte hace % preamble.tex — estilo común de todos los apuntes Froth.
% Cada apunte hace \input{preamble} justo después de \documentclass.
\usepackage[margin=2.4cm]{geometry}
\usepackage{xcolor}
\definecolor{litblue}{RGB}{31,79,121}
\definecolor{litgreen}{RGB}{27,94,32}
\definecolor{litorange}{RGB}{230,81,0}
\usepackage{parskip}
\usepackage{titlesec}
\titleformat{\section}{\Large\bfseries\color{litblue}}{\thesection}{0.6em}{}
\titleformat{\subsection}{\large\bfseries\color{litblue}}{\thesubsection}{0.6em}{}
\usepackage{tcolorbox}
\usepackage{fancyhdr}
\pagestyle{fancy}
\fancyhf{}
\fancyhead[L]{\small\color{litblue}Froth — Apuntes de ML}
\fancyhead[R]{\small\thepage}
\renewcommand{\headrulewidth}{0.4pt}
\usepackage[colorlinks=true,linkcolor=litblue,urlcolor=litblue]{hyperref}

% Cabecera de cada apunte: \cabecera{numero}{titulo}
\newcommand{\cabecera}[2]{%
  \begin{tcolorbox}[colback=litblue,coltext=white,colframe=litblue,arc=2mm]
    {\Large\bfseries Apunte #1 — #2}\\[3pt]
    {\small Proyecto Froth · aprender Machine Learning construyendo}
  \end{tcolorbox}\vspace{0.8em}}

% Cajas temáticas
\newtcolorbox{concepto}[1]{colback=blue!4,colframe=litblue,title={Concepto: #1},arc=1.5mm}
\newtcolorbox{practica}{colback=green!5,colframe=litgreen,title={Pruébalo tú (teclado en mano)},arc=1.5mm}
\newtcolorbox{ojo}{colback=orange!8,colframe=litorange,title={Ojo / error típico},arc=1.5mm}
\newtcolorbox{resumen}{colback=gray!8,colframe=gray!60,title={En una frase},arc=1.5mm}

\newcommand{\codigo}[1]{\texttt{#1}}
 justo después de \documentclass.
\usepackage[margin=2.4cm]{geometry}
\usepackage{xcolor}
\definecolor{litblue}{RGB}{31,79,121}
\definecolor{litgreen}{RGB}{27,94,32}
\definecolor{litorange}{RGB}{230,81,0}
\usepackage{parskip}
\usepackage{titlesec}
\titleformat{\section}{\Large\bfseries\color{litblue}}{\thesection}{0.6em}{}
\titleformat{\subsection}{\large\bfseries\color{litblue}}{\thesubsection}{0.6em}{}
\usepackage{tcolorbox}
\usepackage{fancyhdr}
\pagestyle{fancy}
\fancyhf{}
\fancyhead[L]{\small\color{litblue}Froth — Apuntes de ML}
\fancyhead[R]{\small\thepage}
\renewcommand{\headrulewidth}{0.4pt}
\usepackage[colorlinks=true,linkcolor=litblue,urlcolor=litblue]{hyperref}

% Cabecera de cada apunte: \cabecera{numero}{titulo}
\newcommand{\cabecera}[2]{%
  \begin{tcolorbox}[colback=litblue,coltext=white,colframe=litblue,arc=2mm]
    {\Large\bfseries Apunte #1 — #2}\\[3pt]
    {\small Proyecto Froth · aprender Machine Learning construyendo}
  \end{tcolorbox}\vspace{0.8em}}

% Cajas temáticas
\newtcolorbox{concepto}[1]{colback=blue!4,colframe=litblue,title={Concepto: #1},arc=1.5mm}
\newtcolorbox{practica}{colback=green!5,colframe=litgreen,title={Pruébalo tú (teclado en mano)},arc=1.5mm}
\newtcolorbox{ojo}{colback=orange!8,colframe=litorange,title={Ojo / error típico},arc=1.5mm}
\newtcolorbox{resumen}{colback=gray!8,colframe=gray!60,title={En una frase},arc=1.5mm}

\newcommand{\codigo}[1]{\texttt{#1}}
 justo después de \documentclass.
\usepackage[margin=2.4cm]{geometry}
\usepackage{xcolor}
\definecolor{litblue}{RGB}{31,79,121}
\definecolor{litgreen}{RGB}{27,94,32}
\definecolor{litorange}{RGB}{230,81,0}
\usepackage{parskip}
\usepackage{titlesec}
\titleformat{\section}{\Large\bfseries\color{litblue}}{\thesection}{0.6em}{}
\titleformat{\subsection}{\large\bfseries\color{litblue}}{\thesubsection}{0.6em}{}
\usepackage{tcolorbox}
\usepackage{fancyhdr}
\pagestyle{fancy}
\fancyhf{}
\fancyhead[L]{\small\color{litblue}Froth — Apuntes de ML}
\fancyhead[R]{\small\thepage}
\renewcommand{\headrulewidth}{0.4pt}
\usepackage[colorlinks=true,linkcolor=litblue,urlcolor=litblue]{hyperref}

% Cabecera de cada apunte: \cabecera{numero}{titulo}
\newcommand{\cabecera}[2]{%
  \begin{tcolorbox}[colback=litblue,coltext=white,colframe=litblue,arc=2mm]
    {\Large\bfseries Apunte #1 — #2}\\[3pt]
    {\small Proyecto Froth · aprender Machine Learning construyendo}
  \end{tcolorbox}\vspace{0.8em}}

% Cajas temáticas
\newtcolorbox{concepto}[1]{colback=blue!4,colframe=litblue,title={Concepto: #1},arc=1.5mm}
\newtcolorbox{practica}{colback=green!5,colframe=litgreen,title={Pruébalo tú (teclado en mano)},arc=1.5mm}
\newtcolorbox{ojo}{colback=orange!8,colframe=litorange,title={Ojo / error típico},arc=1.5mm}
\newtcolorbox{resumen}{colback=gray!8,colframe=gray!60,title={En una frase},arc=1.5mm}

\newcommand{\codigo}[1]{\texttt{#1}}

\begin{document}
\cabecera{10}{Cuadernos Jupyter: tu buscador interactivo}

\section{Qué es un notebook y por qué importa}

\begin{concepto}{Jupyter notebook}
Un \textbf{notebook} (archivo \codigo{.ipynb}) es un documento hecho de \textbf{celdas} que
puedes ejecutar una a una. Hay dos tipos de celda: de \textbf{código} (Python que corre) y de
\textbf{texto} (explicaciones en Markdown). Ejecutas una celda con \textbf{Shift + Enter}.
\end{concepto}

La gracia: pruebas, ves el resultado justo debajo, cambias algo y vuelves a correr --- sin
reiniciar todo el programa. Por eso es \textbf{la} herramienta de exploración en data science
y ciencia de datos. Además, un notebook bien hecho es una pieza excelente de portafolio:
cuenta una historia (texto + código + resultados) que cualquiera puede seguir.

\section{Tu cuaderno: \codigo{4\_Notebooks/01\_buscador\_semantico.ipynb}}

Envuelve la función \codigo{most\_similar()} en una interfaz cómoda:
\begin{enumerate}
  \item Una celda de \textbf{preparación} carga los 126 papers y sus vectores.
  \item Una celda de \textbf{consulta} donde cambias la variable \codigo{my\_query} por la
        frase que quieras y ejecutas.
  \item El resultado aparece como \textbf{tabla} (score, título, año, citas), porque
        \codigo{most\_similar} devuelve un DataFrame y Jupyter los dibuja bonitos.
\end{enumerate}

\begin{practica}
Ábrelo tú y valida como experto. En PowerShell, dentro de la carpeta del proyecto:
\begin{enumerate}
  \item \codigo{cd "C:\textbackslash Users\textbackslash Bruce Wayne\textbackslash Documents\textbackslash Master Thesis AI"}
  \item \codigo{.venv\textbackslash Scripts\textbackslash python.exe -m jupyter notebook "4\_Notebooks\textbackslash 01\_buscador\_semantico.ipynb"}
\end{enumerate}
Se abre en tu navegador. Ejecuta las celdas con Shift+Enter, cambia \codigo{my\_query} y
prueba frases de tu campo. \textbf{Tú juzgas} si los papers que salen son los correctos: esa
es la validación que cierra la Fase 2.
\end{practica}

\begin{ojo}
El notebook empieza con \codigo{sys.path.insert(...)} apuntando a la carpeta del proyecto.
Es para que Python encuentre el paquete \codigo{froth} aunque el cuaderno viva en la subcarpeta
\codigo{4\_Notebooks/}. Sin esa línea, \codigo{from froth import embed} fallaría.
\end{ojo}

\section{Notebook vs.\ código (\codigo{.py}): ¿cuándo cada uno?}

\begin{center}
\begin{tabular}{p{6.5cm} p{6.5cm}}
\textbf{Notebook (\codigo{.ipynb})} & \textbf{Módulo (\codigo{.py})} \\
\hline
Explorar, probar ideas, enseñar, presentar resultados. & Código estable que se reutiliza
(el ``motor'', como \codigo{embed.py}). \\
Se lee de arriba a abajo como una historia. & Se importa desde otros archivos. \\
\end{tabular}
\end{center}
Regla sana: exploras en el notebook; cuando algo funciona y lo vas a reutilizar, lo mueves a
un \codigo{.py} del paquete. Tu buscador ya vive en \codigo{embed.py}; el notebook solo lo
\emph{usa}.

\begin{resumen}
Un notebook es un documento de celdas ejecutables (Shift+Enter) ideal para explorar y
presentar. El tuyo envuelve \codigo{most\_similar} para que escribas consultas y veas los
papers en tabla: úsalo para validar la Fase 2 con tu criterio de experto.
\end{resumen}

\end{document}
